\paragraph{Temporal control (development level).} Because the FEniCSx step ($0.02\,\mathrm{s}$) is ten times the solids4foam step, the Gc level was rerun from $t=0$ with $\Delta t=0.005\,\mathrm{s}$ (same mesh and numerics, preregistered question and criterion in \texttt{prereg/ADDENDUM\_A7.md}). The tip changes by at most 7.0\,\% (rms 3.1\,\%) for $t\ge0.2\,\mathrm{s}$ (and by 17\,\% at $t=0.1\,\mathrm{s}$, where the tip is micrometres), and the peak of $F_y$ near $0.67\,\mathrm{s}$ rises from 237 to 472\,$\mu$N, against 563\,$\mu$N in solids4foam (72\,\% of the gap closed): (sample maxima: the $\Delta t=0.02\,\mathrm{s}$ peak is under-sampled, so this fraction is uncertain); on the accepted level Gmi the sample maximum of $F_y$ is only 290\,$\mu$N at $t=0.69\,\mathrm{s}$, about half of the solids4foam peak, and $F_y$ has the opposite sign at $0.7\,\mathrm{s}$ (Table~\ref{tab:matched}). The height of the $F_y$ peak is therefore largely a time-step effect in the FEniCSx data at $\Delta t=0.02\,\mathrm{s}$ (shown on the development level only), and the tip agreement of Table~\ref{tab:matched} contains temporal errors of the same order as the differences it reports (at $\Delta t=0.005\,\mathrm{s}$ the development-level tip is $-6\,\%$ below solids4foam at $0.6\,\mathrm{s}$ and $+2\,\%$ above at $0.75\,\mathrm{s}$). This is one mesh and two time steps: it is neither a temporal convergence study nor an error bound. The step used in Table~\ref{tab:matched} was not changed to improve agreement.\par\begin{table}[htbp]\centering\caption{Temporal control on the development level Gc (not an accepted level) against solids4foam (ours).}\label{tab:temporal}\scriptsize\setlength{\tabcolsep}{3pt}\temporaltable\end{table}
